% 340_Neutrino_Galois_De_ch.tex
% Johann Pascher, 30. Aug. 2026

\providecommand{\BM}{\textbf{[B]}}
\providecommand{\KM}{\textbf{[K]}}
\providecommand{\SM}{\textbf{[S]}}
\providecommand{\SetzM}{\textbf{[SETZUNG]}}

\chapter*{Neutrino-Massenhierarchie aus der $GF(27)^*$-Frobenius-Struktur}
\addcontentsline{toc}{chapter}{Neutrino-Massenhierarchie aus $GF(27)^*$}

\begin{abstract}
Ausgehend von der Frobenius-Zerlegung $GF(27)^*\cong\mathbb{Z}_{26}\cong
\mathbb{Z}_2\times\mathbb{Z}_{13}$ (Dok.~339) werden die gemessenen
Neutrino-Massendifferenzen $\Delta m^2_\text{atm}$ und $\Delta m^2_\text{sol}$
aus algebraischen Eigenschaften der vier Frobenius-Orbits in $\mathbb{Z}_{13}$
abgeleitet. Die atmosphärische Massendifferenz
$\Delta m^2_\text{atm}=(11^2-1)\,m_\nu^2=120\,m_\nu^2$ (1.0\,\% Abweichung)
und die solare $\Delta m^2_\text{sol}=(|\mathbb{Z}_{13}|-2)/3\cdot m_\nu^2=(11/3)\,m_\nu^2$
(0.46\,\% Abweichung) folgen aus der Zahl~$11$, dem maximalen Element des
vierten Frobenius-Orbits $\{7,8,11\}$ in $\mathbb{Z}_{13}$. Orbit~3
$\{4,10,12\}$ ist algebraisch selbstinvers und trägt Majorana-Charakter.
Die kohärente Venting-Amplitude $p_\text{vent}=5/24$ wird aus der
Branching-Rule $\rho_4\!\downarrow_{A_4}=\rho_3^{A_4}\oplus\rho_1^{A_4}$
und dem Index $|A_5:A_4|=5$ algebraisch bewiesen. Die SICC-Venting-Formel
ist äquivalent zur Standard-Oszillationsformel. Status: \BM\ für die
Gruppentheorie; \KM\ für die Massendifferenzen.
\end{abstract}

% -------------------------------------------------------------------
\section{Ausgangspunkt}

In Dok.~339 wird die Frobenius-Zerlegung
\[
GF(27)^*\cong\mathbb{Z}_{26}\cong\mathbb{Z}_2\times\mathbb{Z}_{13}
\]
gesichert \BM. Die vier Dreier-Orbits von $\mathbb{Z}_{13}$ unter dem
Frobenius $\phi:x\mapsto3x$ sind:

\begin{center}
\begin{tabular}{clrr}
\toprule
Orbit & Elemente & Summe & Produkt mod~13 \\
\midrule
1 & $\{1,3,9\}$   & 13 & 1  \\
2 & $\{2,5,6\}$   & 13 & 8  \\
3 & $\{4,10,12\}$ & 26 & 12 \\
4 & $\{7,8,11\}$  & 26 & 5  \\
\bottomrule
\end{tabular}
\end{center}

Die FFGFT-Neutrinobasis-Masse ist spekulativ als
$m_\nu=\xi^2/2\cdot m_e\approx4.54\,\text{meV}$ angesetzt \SM\
(Dok.~047). Die vorliegenden Rechnungen testen, ob die gemessenen
Massendifferenzen aus der Galois-Struktur folgen.

% -------------------------------------------------------------------
\section{Algebraische Eigenschaften der Orbits \BM}

\subsection{Orbit~4 als Vielfaches von Orbit~1}

\[
\{7,8,11\}=7\cdot\{1,3,9\}\pmod{13},
\]
denn $7\cdot1=7$, $7\cdot3=21\equiv8$, $7\cdot9=63\equiv11$ (mod~13).
Da $2^{-1}\equiv7\pmod{13}$:
\[
\boxed{\text{Orbit}_4=2^{-1}\cdot\text{Orbit}_1\pmod{13}.}
\]

\subsection{Dualität Orbit~2 $\leftrightarrow$ Orbit~4}

Die minimalen Elemente und ihre multiplikativen Inversen in $\mathbb{Z}_{13}$:

\begin{center}
\begin{tabular}{cccc}
\toprule
Orbit & min & $\text{min}^{-1}\!\pmod{13}$ & liegt in \\
\midrule
1 & 1 & 1 & Orbit~1 \\
2 & 2 & 7 & Orbit~4 \\
3 & 4 & 10 & Orbit~3 \\
4 & 7 & 2 & Orbit~2 \\
\bottomrule
\end{tabular}
\end{center}

Orbits~2 und~4 sind algebraisch invers: $2^{-1}=7$, $7^{-1}=2$.
Orbit~3 ist selbstinvers: $4^{-1}=10$, $10^{-1}=4$, $12^{-1}=12$
(mod~13). Die Selbstinversheit von Orbit~3 entspricht algebraisch
einem Majorana-Charakter. Die Summe:
\[
4+10+12=26=|GF(27)^*|.
\]

% -------------------------------------------------------------------
\section{Neutrinomassen aus der Orbit-Struktur \KM}

\subsection{Der Faktor~11}

Das maximale Element des vierten Orbits, $\max(\{7,8,11\})=11$, folgt aus:
\[
11=|\mathbb{Z}_{13}|-|\text{Frobenius-Fixpunkte in }\mathbb{Z}_{13}|-1
  =13-1-1=11.
\]
Dieses Ergebnis bestimmt die Massenhierarchie (normale Ordnung):

\begin{center}
\begin{tabular}{clll}
\toprule
Zustand & Galois-Herkunft & Masse & Wert \\
\midrule
$\nu_1$ & Orbit~1 (D4, 4D) & $m_\nu$ & $4.54\,\text{meV}$ \\
$\nu_2$ & Orbit~3 (selbstinvers, E6, solar) & $\sqrt{14/3}\,m_\nu$ & $9.81\,\text{meV}$ \\
$\nu_3$ & Orbit~4 (E8, atm., schwerste) & $11\,m_\nu$ & $49.96\,\text{meV}$ \\
\bottomrule
\end{tabular}
\end{center}

Die Summe $\sum m_i=(1+\sqrt{14/3}+11)\,m_\nu\approx64.3\,\text{meV}$
ist konsistent mit der Planck-2018-Grenze $\sum m_i<120\,\text{meV}$.

\subsection{Atmosphärische Massendifferenz}

\[
\Delta m^2_\text{atm}=m_{\nu_3}^2-m_{\nu_1}^2
=(11^2-1)\,m_\nu^2=120\,m_\nu^2
=2.476\times10^{-3}\,\text{eV}^2.
\]
Gemessen: $2.500\times10^{-3}\,\text{eV}^2$.
\textbf{Abweichung: 1.0\,\%.}

Der Faktor $120=|A_5\times\mathbb{Z}_2|=11^2-1$ ist die Ordnung der
vollen ikosaedrischen Symmetriegruppe -- konsistent mit der
$A_5$-Symmetrie des HLV-Sektors (Dok.~285).

\subsection{Solare Massendifferenz}

\[
\Delta m^2_\text{sol}=m_{\nu_2}^2-m_{\nu_1}^2
=\left(\frac{14}{3}-1\right)m_\nu^2=\frac{11}{3}\,m_\nu^2
=7.565\times10^{-5}\,\text{eV}^2.
\]
Gemessen: $7.530\times10^{-5}\,\text{eV}^2$.
\textbf{Abweichung: 0.46\,\%.}

Die algebraische Herleitung des Faktors $11/3$ \BM:
\[
\frac{11}{3}
=\frac{|\mathbb{Z}_{13}|-|\text{Fixpunkte von }GF(27)^*|}{|\text{Frobenius-Ordnung}|}
=\frac{13-2}{3}.
\]
Zähler: die zwei Fixpunkte $\{+1,-1\}$ von $GF(27)^*$ unter dem
Frobenius werden von $|\mathbb{Z}_{13}|=13$ subtrahiert.
Nenner: die Ordnung des Frobenius $\phi:x\mapsto x^3$ ist~3.

% -------------------------------------------------------------------
\section{Kappa und die SICC-Verbindung \KM}

\subsection{Verbindung $\kappa$ -- $\Delta m^2_\text{atm}$}

Aus Dok.~339 (und der SICC-Ableitung von Paul Masic) gilt:
\[
\kappa=\frac{8}{r_e}\cdot m_e\,\xi^2=6\,\xi^2\,m_e=12\,m_\nu.
\]
Damit folgt direkt:
\[
\Delta m^2_\text{atm}=(11^2-1)\,m_\nu^2=120\,m_\nu^2
=\frac{5}{6}\,\kappa^2.
\]
Das Verhältnis:
\[
\frac{\kappa^2}{\Delta m^2_\text{atm}}=\frac{144}{120}=\frac{6}{5}.
\]

\subsection{Die kohärente Venting-Amplitude $p_\text{vent}=5/24$ \BM}

Wenn die W-Venting-Ereignisse nicht als Energieverlust sondern als
kohärente Quantenphasen akkumuliert werden, lautet die
Oszillationsphase:
\[
\phi_\text{venting}=p_\text{vent}\cdot\frac{\kappa^2\,L}{\hbar c\,E_\nu}.
\]
Damit $\phi_\text{venting}=\Delta m^2_\text{atm}\cdot L/(\hbar c\,E_\nu)$,
muss gelten:
\[
p_\text{vent}=\frac{\Delta m^2_\text{atm}}{4\,\kappa^2}
=\frac{120\,m_\nu^2}{4\cdot144\,m_\nu^2}=\frac{120}{576}=\frac{5}{24}.
\]

\subsection{Algebraische Herleitung von $p=5/24$ aus der Gruppenstruktur \BM}

Die Branching-Rule für die Einschränkung der 4D-Darstellung
$\rho_4$ von $A_5$ auf die Untergruppe $A_4$ (Tetraedergruppe,
Ordnung~12) lautet:
\[
\boxed{\rho_4\!\downarrow_{A_4}=\rho_3^{A_4}\oplus\rho_1^{A_4}.}
\]
Diese folgt aus der $A_4$-Charaktertafel: die Charaktere von
$\rho_4$ auf den $A_4$-Klassen $(e,\,(123),\,(132),\,(12)(34))$
sind $(4,1,1,0)$, und das Skalarprodukt mit den
$A_4$-Charakteren ergibt Multiplizitäten $1+0+0+0=1$ für
$\rho_1^{A_4}$ und $1$ für $\rho_3^{A_4}$.

\medskip
\textbf{Konsequenz:} $A_4$ lässt in $\rho_4$ eine
eindimensionale Unterdarstellung (die W-Richtung) fest.
Das heißt, $A_4$ ist der \emph{Stabilisator} der W-Richtung
in $A_5$.

In $A_5$ gibt es genau
\[
|A_5:A_4|=\frac{60}{12}=5
\]
verschiedene Konjugate von $A_4$ -- entsprechend den 5~Tetraedern
die im Ikosaeder eingebettet sind. Jedes Konjugat stabilisiert
eine andere W-Richtung. Das W-Venting wählt genau eine dieser
5~Richtungen (die physikalische 4.~Raumrichtung), daher:
\[
\boxed{p_\text{vent}=\frac{|A_5:A_4|}{\text{Kissing}(D_4)}
=\frac{5}{24}.}
\]

\subsection{Äquivalenz zur Standard-Oszillationsformel}

Die vollständige SICC-Venting-Formel lautet:
\[
\phi_\text{venting}
=\frac{5}{24}\cdot\frac{\kappa^2\,L}{\hbar c\,E_\nu}
=\frac{5}{24}\cdot\frac{144\,m_\nu^2\,L}{\hbar c\,E_\nu}
=\frac{120\,m_\nu^2\cdot L}{4\,\hbar c\,E_\nu}
=\frac{\Delta m^2_\text{atm}\cdot L}{4\,\hbar c\,E_\nu}.
\]
Das ist \emph{identisch} mit der Standard-Oszillationsformel.
Alle Parameter -- $\kappa$, $p_\text{vent}$, $\Delta m^2_\text{atm}$
-- kommen aus der Galois-Struktur $GF(27)^*$ und der
Gruppenstruktur $A_5\supset A_4$. Kein freier Parameter,
keine Kalibrierung an DUNE.

Die vollständige Ableitungskette:
\[
GF(27)^*
\xrightarrow{[B]}\;m_\nu
\xrightarrow{[K]}\;\kappa=12\,m_\nu
\xrightarrow{[K]}\;\Delta m^2_\text{atm}=120\,m_\nu^2
\xrightarrow{[B]}\;p_\text{vent}=\frac{|A_5:A_4|}{\text{Kissing}(D_4)}
=\frac{5}{24}.
\]

\subsection{Nichtschließung der Winkelstruktur: $\mathbb{Z}_{13}\times\mathbb{Z}_{24}$ \BM}

Die SICC-Simulation quantisiert Richtungen auf die 24
Kissing-Vektoren von $D_4$; die Frobenius-Struktur liefert
Orbit-Phasen $2\pi k/13$. Beide Winkelgitter sind teilerfremd:
\[
\gcd(13,24)=1,\qquad \operatorname{kgV}(13,24)=312.
\]
Die kombinierte Winkelstruktur schließt daher erst nach 312
Schritten, nicht nach 24 -- eine diskrete Spirale. Das ist die
Nichtabschluss-Struktur von Dok.~193 in Winkelform: dort auf der
Ebene der Massenprodukte $r_i\cdot r_j$, hier auf der Ebene
$\mathbb{Z}_{13}\times\mathbb{Z}_{24}$.

Kontinuierlich tritt die Spirale zusätzlich über das Beam-Spektrum
auf: $\phi(E)\propto1/E$ ist nicht periodisch in $E$. Über
$[0.5,\,5]\,\text{GeV}$ (DUNE) windet sich die Phase um
$1.155$ Umläufe mit Rest $55.7°$ mod $2\pi$ -- die Kurve schließt
nicht. Im 1D-Phasenhistogramm ist das unsichtbar; im Polarplot
(Radius $=$ Spektralgewicht, Winkel $=\phi\bmod2\pi$) wird die
Spirale sichtbar.

\subsection{Fraktale Korrektur: nicht anwendbar auf Verhältnisse}

Die fraktale Korrektur $K_\text{frak}=1-100\,\xi$ (Dok.~011/133)
betrifft absolute Massen, nicht dimensionslose Verhältnisse
(Dok.~306). $\Delta m^2_\text{atm}/m_\nu^2=120$ ist ein Verhältnis;
$K_\text{frak}$ darf hier nicht angewendet werden. Zur Kontrolle:
angewendet würde die Abweichung von $\Delta m^2_\text{atm}$ von
$1.0\,\%$ auf $2.3\,\%$ steigen; der Effekt auf die
DUNE-Verschwindung wäre $0.2$ Prozentpunkte -- unter Detektorauflösung.

% -------------------------------------------------------------------
\section{Mischungswinkel \KM\,/\,\SM}

Zwei Mischungswinkel lassen sich mit $<6\,\%$ Abweichung aus
Orbit-2-Elementen ablesen \KM:
\[
\sin^2\theta_{12}\approx\cos^2\!\left(\frac{2\pi\cdot2}{13}\right)=0.323
\quad(\text{gemessen: }0.307,\;\text{Abw.\ }5.1\,\%)
\]
\[
\sin^2\theta_{23}\approx\cos^2\!\left(\frac{2\pi\cdot5}{13}\right)=0.560
\quad(\text{gemessen: }0.545,\;\text{Abw.\ }2.8\,\%)
\]

Der Reaktorwinkel $\theta_{13}$ hat zwei unabhängige Galois-Herleitungen \KM:

\textbf{(a) Galois-Formel aus Dok.~328 (stärker):}
\[
\theta_{13}=\arcsin\!\left((25\,\xi)^{1/3}\right),\quad
25\,\xi=\frac{1}{300}=\frac{1}{5\cdot|A_5|},
\]
\[
\sin^2\theta_{13}=\left(\frac{1}{300}\right)^{2/3}\approx0.02230
\quad(\text{gemessen: }0.0220,\;\text{Abw.\ }1.4\,\%)
\]
Der Faktor $25=5^2$ entstammt der Galois-Identität
$43200=|GF(9)^*|^2\cdot5^2\cdot|GF(27)|=64\cdot25\cdot27$
(Dok.~338), wobei $5^2$ das Quadrat der maximalen $A_5$-Darstellungsdimension ist.

\textbf{(b) Summenregel aus $R_\nu$ (unabhängige Bestätigung):}
\[
\sin^2\theta_{13}=\frac{2}{3}\,R_\nu
=\frac{2}{3}\cdot\frac{11}{360}=\frac{11}{540}\approx0.02037
\quad(\text{Abw.\ }7.4\,\%)\quad\cite{Razzaghi2022}
\]
In Galois-Lesart: $2/3=\min(\text{Orbit}_2)/\max(\text{Orbit}_1)$.

Dieselbe Quelle liefert unabhängig:
\[
\sin^2\theta_{12}=\frac{1}{3-2R_\nu}\approx0.340
\quad(\text{Abw.\ zu Galois-Wert }0.323{:}\ 5.4\,\%)
\]

Die CP-Verletzungsphase $\delta_{CP}\approx-\pi/2$ (experimenteller
Hinweis) entspricht dem Orbit-3-Element $k=10$:
$\arg(e^{2\pi i\cdot10/13})=-83.1°$, Abweichung $6.9°$ \SM.

\subsection*{Literaturverknüpfung: T$_{13}$ als Flavour-Symmetriegruppe}

Die Frobenius-Gruppe $T_{13}=\mathbb{Z}_{13}\rtimes\mathbb{Z}_3$
ist in der Literatur als Kandidat für die Flavour-Symmetrie der
Neutrinos bekannt \cite{Hartmann2011}. Sie ist identisch mit der
aus $GF(27)^*$ abgeleiteten Frobenius-Struktur dieses Dokuments.
Die Literatur behandelte $T_{13}$ als abstrakte Symmetriegruppe
(ohne physikalische Herleitung); FFGFT liefert den algebraischen
Ursprung: $T_{13}=GF(27)^*$ unter dem Frobenius-Automorphismus
$\phi:x\mapsto x^3$.

% -------------------------------------------------------------------
\section{Orbit~3 als Majorana-Sektor \BM}

Orbit~3 $\{4,10,12\}$ ist vollständig selbstinvers in
$\mathbb{Z}_{13}$:
\[
4^{-1}\equiv10,\quad 10^{-1}\equiv4,\quad 12^{-1}\equiv12\pmod{13}.
\]
Die Selbstinversheit bedeutet: die zu Orbit~3 gehörigen
Teilchen sind ihre eigenen Antiteilchen -- Majorana-Charakter.
Zusätzlich gilt $4+10+12=26=|GF(27)^*|$ \BM.

Falls sterile Neutrinos existieren, liegen sie algebraisch
im Orbit-3-Sektor mit minimaler Masse
$m_\text{steril,min}=4\,m_\nu=18.2\,\text{meV}$ \SM.

% -------------------------------------------------------------------
\section{Neutrinoloser Doppelbeta-Zerfall \KM}

Mit den gemessenen PMNS-Mischungswinkeln und den Galois-Massen:
\[
m_{ee}=\left|c_{12}^2c_{13}^2\,m_1+s_{12}^2c_{13}^2\,m_2
+s_{13}^2\,m_3\right|\approx7.1\,\text{meV}\quad(\text{Majorana-Phasen null}).
\]
Die KamLAND-Zen-Grenze $m_{ee}<36\,\text{meV}$ wird deutlich
eingehalten. Der Bereich bei destruktiver Interferenz:
$m_{ee,\min}\approx1.2\,\text{meV}$.

% -------------------------------------------------------------------
\section{Falsifizierbare Vorhersagen}

\begin{enumerate}
\item \textbf{Neutrinomassen (normale Hierarchie):}
  $m_1=4.54$, $m_2=9.81$, $m_3=49.96\,\text{meV}$,
  $\sum=64.3\,\text{meV}$.
  Testbar: CMB-S4, Euclid (Kosmologie auf 1\,meV-Niveau).
  Unabhängige Konvergenz: Razzaghi et al.\ \cite{Razzaghi2022}
  extrahieren aus A$_4$-Textur und Experiment
  $m_2=(9.21$--$9.92)$, $m_3=(49.91$--$51.42)$,
  $\sum=(63.97$--$64.55)\,\text{meV}$; unsere Galois-Werte
  $m_2, m_3$ und $\sum$ liegen innerhalb dieser Bereiche.

\item \textbf{Massendifferenzen-Ratio} \KM:
  \[\frac{\Delta m^2_\text{atm}}{\Delta m^2_\text{sol}}
  =\frac{3(11^2-1)}{11}=32.73\quad(\text{gemessen: }33.20).\]
  Sobald $m_\nu$ auf 1\,meV bekannt ist, ist $\Delta m^2/m_\nu^2$
  direkt testbar.

\item \textbf{Orbit-3-Sektor (Majorana)} \KM:
  Sterile Neutrinos (falls existent) sind Majorana-Teilchen mit
  $m_\text{steril,min}=4\,m_\nu=18.2\,\text{meV}$.
  Testbar: IceCube, SBN (Fermilab), BEST-Experiment.

\item \textbf{Venting-Rate und DUNE-Observable} \KM:
  Die kohärente SICC-Venting-Amplitude
  $p_\text{vent}=5/24\approx20.8\,\%$ pro Kohärenzschritt ist
  eine modell-interne Vorhersage. Die DUNE-Observable hängt vom
  Beam-Spektrum ab: monochromatisch $2.5\,\text{GeV}$ ergibt
  $99.8\,\%$ $\nu_\mu$-Verschwindung (erstes Oszillationsmaximum,
  $\phi=1.61\,\text{rad}$); ein gauß\-förmiges Spektrum
  $2.5\pm1.0\,\text{GeV}$, begrenzt auf $[0.5,\,5]\,\text{GeV}$,
  ergibt gemittelt $76\,\%$ (SICC~V3: $76.16\,\%$). Für die
  direkte Vorhersage ist die reale LBNF-Flusstabelle als
  Gewicht einzusetzen.
\end{enumerate}

% -------------------------------------------------------------------
\section{Konjugationsstruktur der Orbits \BM}

Die vier Frobenius-Orbits in $\mathbb{Z}_{13}$ bilden unter der
Konjugation $x\mapsto-x\pmod{13}$ genau zwei konjugierte
3/\={3}-Paare -- identisch mit den Konjugationsklassen $C_{3_1},
C_{\bar3_1}, C_{3_2}, C_{\bar3_2}$ von $T_{13}$
\cite{Hartmann2011}:
\[
-\text{Orbit}_1=\text{Orbit}_3,\quad
-\text{Orbit}_2=\text{Orbit}_4.
\]
Zusätzlich zur Inversions-Dualität ($x\mapsto x^{-1}$, Dok.~339)
erfüllen die Orbits damit eine vollständige
$\mathbb{Z}_2\times\mathbb{Z}_3$-Symmetrie in $\mathbb{Z}_{13}$.
Orbit~1 ist auch unter Inversion selbstinvers ($3^{-1}=9$,
$9^{-1}=3$, beide in Orbit~1).

\section{Modell-Diskrepanzen zur Einordnung \SM}

Das A$_4$/TM$_2$-Modell \cite{Razzaghi2022}, das die
$\theta_{13}$-Summenregel liefert, trifft folgende Vorhersagen
die vom Galois-Bild und vom Experiment abweichen:
\begin{itemize}
\item $\sin^2\theta_{23}=0.433$--$0.441$ (unteres Oktant) --
  gemessen: $0.545$--$0.560$ (oberes Oktant); Galois: $0.560$ (oberes Oktant).
\item $\delta_{CP}\approx\pm51°$ --
  T2K-Hinweis: $\approx-90°$; Galois-Kandidat ($k=10$ aus Orbit~3,
  Begründung offen \SM): $\approx-83°$.
\item $m_{ee}=5.1$--$5.3\,\text{meV}$ (mit Majorana-Phasen) --
  Galois (Phasen null): $7.1\,\text{meV}$.
\end{itemize}
In $\theta_{23}$ und $\delta_{CP}$ stimmen die
FFGFT-Galois-Werte besser mit dem Experiment überein als das
A$_4$-Modell. Die Summenregel für $\theta_{13}$ wird übernommen,
die übrigen A$_4$-Vorhersagen nicht.

\section{Offene Punkte}

\begin{enumerate}
\item Verbindung der diskreten Orbit-Phasen $2\pi k/13$ mit der
  PMNS-$U$-Matrix \SM.
\item Algebraische Begründung warum $k=10$ aus Orbit~3 für $\delta_{CP}$
  ausgezeichnet ist; Kandidat $\arg(e^{2\pi i\cdot10/13})=-83{,}1°$
  (T2K-Hinweis $\approx-90°$, Abw.\ $6{,}9°$) \SM.
\end{enumerate}

% -------------------------------------------------------------------
\section{Entstehungshinweis}

Dieses Dokument wurde am 30.~August 2026 in Zusammenarbeit mit dem
KI-Assistenten Claude (Anthropic, Sonnet~4.6) erarbeitet. Die algebraischen
Rechnungen (Frobenius-Orbits in $\mathbb{Z}_{13}$, Branching-Rule
$\rho_4\!\downarrow_{A_4}$, Charaktertheorie von $A_5$ und $A_4$,
numerische Prüfung der Massendifferenzen) wurden vom Assistenten
durchgeführt und vom Autor geprüft. Die physikalische Interpretation,
die Verbindung zum SICC-Rahmen und alle Statusmarkierungen
(\BM, \KM, \SM) liegen beim Autor.

% -------------------------------------------------------------------
\section{Zeichenerklärung}

\begin{tabular}{ll}
\textbf{[SETZUNG]} & Axiom oder deklarierte externe Eingabe -- nicht hergeleitet \\
\textbf{[K]}       & Kern -- aus $\xi$ und den Setzungen hergeleitet, numerisch geprüft \\
\textbf{[B]}       & Brücke -- algebraisch bewiesen \\
\textbf{[S]}       & Skizze -- plausibel, noch nicht vollständig ausgeführt \\
\end{tabular}

% -------------------------------------------------------------------
\section{Prüfskripte}

\texttt{pruef\_340\_neutrino\_galois.py} (13/13 Assertions bestanden);\\
\texttt{pruef\_340b\_spirale\_breitband.py} (6/6 Assertions bestanden:
Nichtschließung $\mathbb{Z}_{13}\times\mathbb{Z}_{24}$, Spirale,
Breitband-DUNE, $K_\text{frak}$).

\emph{Erklärt 30.~August 2026, ergänzt 31.~August 2026 ($\theta_{13}$, T$_{13}$-Literatur, Nichtschließung, Breitband)}

\begin{thebibliography}{9}
\bibitem{Razzaghi2022}
N.~Razzaghi, S.~M.~M.~Rasouli, P.~Parada, P.~V.~Moniz,
\emph{Two-zero textures based on $A_4$ symmetry and unimodular mixing matrix},
arXiv:2203.05753 [hep-ph] (2022).

\bibitem{Hartmann2011}
C.~Hartmann, A.~Zee,
\emph{Neutrino mixing and the Frobenius group $T_{13}$},
arXiv:1106.0333 [hep-ph] (2011).
\end{thebibliography}

